\documentclass[twocolumn,11pt,a4paper]{article}

\usepackage[utf8]{inputenc}
\usepackage[T1]{fontenc}
\usepackage{lmodern}
\usepackage[a4paper,margin=2cm,columnsep=0.6cm]{geometry}
\usepackage{amsmath,amssymb,amsthm,mathtools}
\usepackage{bm}
\usepackage{graphicx}
\usepackage{enumitem}
\usepackage{booktabs}
\usepackage{array}
\usepackage{microtype}
\usepackage{xcolor}
\usepackage[colorlinks=true,linkcolor=blue!45!black,citecolor=blue!45!black,
            urlcolor=blue!45!black]{hyperref}
\usepackage[round]{natbib}

% ---- theorem environments ----
\newtheoremstyle{main}{6pt}{6pt}{\itshape}{}{\bfseries}{.}{.5em}{}
\theoremstyle{main}
\newtheorem{theorem}{Theorem}[section]
\newtheorem{lemma}[theorem]{Lemma}
\newtheorem{proposition}[theorem]{Proposition}
\newtheorem{corollary}[theorem]{Corollary}

\theoremstyle{definition}
\newtheorem{definition}[theorem]{Definition}
\newtheorem{axiom}{Axiom}
\newtheorem{example}[theorem]{Example}

\theoremstyle{remark}
\newtheorem{remark}[theorem]{Remark}

% ---- macros ----
\newcommand{\Real}{\mathbb{R}}
\newcommand{\Nat}{\mathbb{N}}
\newcommand{\Out}{\mathcal{X}}
\newcommand{\Know}{\mathcal{K}}
\newcommand{\Acts}{\mathcal{Y}}
\newcommand{\Cell}{\mathcal{C}}
\newcommand{\Dec}{D}
\newcommand{\Proj}{\Pi}
\newcommand{\Recv}{R}
\newcommand{\Sval}{S}
\newcommand{\floor}{\beta}
\newcommand{\tol}{\tau}
\newcommand{\shift}{\Delta}
\newcommand{\Mean}{\mathcal{M}}
\newcommand{\dist}{d}
\newcommand{\Sig}{\Sigma}
\newcommand{\Traj}{T}
\newcommand{\Comp}{\mathcal{C}}
\newcommand{\Lab}{\mathcal{L}}
\newcommand{\rate}{r}
\newcommand{\window}{\omega}
\newcommand{\thresh}{\theta^*}
\newcommand{\Carrier}{\mathsf{c}}
\newcommand{\pot}{\kappa}
\newcommand{\half}{\tfrac{1}{2}}
\newcommand{\Mon}{\mathcal{M}}        % monitoring receiver
\newcommand{\Tel}{\mathcal{T}}        % telemetry event sequence
\newcommand{\Dem}{D}                  % cross-demand
\newcommand{\quies}{\mathsf{Q}}       % quiescence predicate
\newcommand{\flo}{\beta_{\mathrm{cut}}} % graph floor

\title{\bfseries Zero Decoder-Shift Under Extreme Temporal Compression:\\
A Graph-Theoretic Account of Propagation Equilibrium\\
in Bounded Receivers}

\author{
Kundai Farai Sachikonye\\
Technical University of Munich\\
AIMe Registry for Artificial Intelligence\\
\texttt{kundai.sachikonye@wzw.tum.de}
}
\date{\today}

\begin{document}
\maketitle

% ======================================================================
\begin{abstract}
\noindent
We study the information transfer properties of a finite-duration signal
transmitted through a bounded cognitive receiver at variable playback
rate. We define the \emph{decoder-shift} as the net movement a received
signal induces in a receiver's internal meaning-representation, and prove
that a strictly positive \emph{resolution floor} $\floor > 0$ exists for
every bounded receiver: no signal, however constructed, can drive
semantic distance to zero. We then introduce the \emph{temporal
compression rate} $\rate$ and prove the \emph{Zero Decoder-Shift
Theorem}: for every bounded receiver there exists a critical rate
$\rate^* < \infty$ above which the transmitted signal induces exactly
zero decoder-shift, regardless of its semantic content. We derive
$\rate^*$ in closed form as a function of the receiver's perceptual
integration window $\window$ and the signal duration $\tau$. We further
develop a graph-theoretic account of signal trajectories under
compression using labeled integer compositions, proving that a signal
compressed above $\rate^*$ remains a \emph{complete} trajectory in the
sender's representation graph while being a \emph{null trajectory} in
the receiver's decoder graph. This separation --- completeness on the
sender side, nullity on the receiver side --- constitutes what we call
\emph{propagation equilibrium}: the state in which a signal has been
fully transmitted but no information has propagated across the
sender--receiver boundary. We characterize propagation equilibrium
geometrically, prove it is stable under perturbation of compression
rate, and derive the minimum compression ratio required to guarantee it
for an arbitrary finite-duration signal. We then extend the framework
to a second class of bounded receiver: an \emph{automated monitoring
system} that observes not the signal content but the \emph{telemetry
event sequence} the signal's transmission produces. Applying the theory
of semantic quiescence to this second receiver, we prove the
\emph{Adversarial Quiescence Theorem}: a signal transmitted at rate
$\rate > \rate^*$ produces an authentic telemetry sequence that is
semantically quiescent with the monitoring system's learned model of a
valid transmission, because the events are real, complete, and in the
correct order. No information crosses the human perceptual boundary, yet
every event the monitoring receiver expects to observe is present and
unmodified. The framework applies to any bounded oscillatory transmission
system and has implications for the theory of channel capacity under
perceptual constraints.
\end{abstract}
% ======================================================================

\section{Introduction}
\label{sec:intro}

\subsection{The question}

When a finite-duration signal is transmitted from a sender to a receiver,
under what conditions does the receiver undergo no net change in internal
state? The question is standard in information theory when framed in
terms of channel noise and coding rate
\citep{shannon1948,cover2006}. It is less studied when the limiting
factor is not the channel but the receiver's own finite temporal
resolution: the minimum duration over which the receiver can integrate
an incoming signal into a representation.

We approach this question from the side of the \emph{sender}. Given a
signal of fixed semantic content and fixed duration $\tau$, and given
that the sender controls the rate at which the signal is delivered, we
ask: does there exist a delivery rate above which the receiver's internal
state is unchanged, and if so, what is that rate?

The answer, which we prove as a theorem, is yes. The critical rate
depends on the receiver's perceptual integration window and not on the
signal's content. Above this rate, the signal is \emph{complete} from
the sender's perspective --- every component was transmitted --- yet
\emph{null} from the receiver's perspective --- no decoder-shift
occurred. We call this condition \emph{propagation equilibrium}.

\subsection{Background and motivation}

The notion of a bounded receiver has classical antecedents in bounded
rationality \citep{simon1955,simon1956} and in psychophysics, where
temporal integration windows for visual and auditory stimuli are
well-established empirical quantities \citep{fechner1860,efron1970,
potter1975,thorpe1996}. What is less well-studied is the algebraic
structure of signal transmission under compression rates that exceed
these windows.

The graph-theoretic framework we develop is motivated by the theory of
integer compositions \citep{andrews1976,heubach2009}. A signal of
duration $\tau$ divided into $n$ temporal segments admits $2^{n-1}$
ordered decompositions (compositions), each representing a different
trajectory through the signal's state space. When each segment is
further labeled by one of $d$ semantic dimensions, the total number of
distinct labeled trajectories grows as $d \cdot (d+1)^{n-1}$ --- an
exponential in $n$ that we call \emph{composition-inflation}. The key
structural result is that this exponential growth is a property of the
\emph{sender's representation} of the signal, and is preserved under
temporal compression. The receiver's representation, however, collapses
to a single null trajectory when compression exceeds the perceptual
threshold.

\subsection{Contributions}

\begin{enumerate}[leftmargin=1.4em,itemsep=2pt]
\item \textbf{The bounded receiver and the Floor Theorem}
  (\S\ref{sec:receiver}). We define the bounded receiver formally and
  prove the resolution floor $\floor > 0$ is structurally unavoidable
  for any finite cognitive system.
\item \textbf{Temporal resolution and the integration window}
  (\S\ref{sec:window}). We introduce the perceptual integration window
  $\window$ and prove that signals shorter than $\window$ are
  decoder-invisible: they induce zero shift.
\item \textbf{The Zero Decoder-Shift Theorem}
  (\S\ref{sec:zero}). We prove that for every signal of duration $\tau$
  and every bounded receiver with window $\window$, the critical
  compression rate $\rate^* = \tau / \window$ exists and is finite, and
  that playback at any rate $\rate > \rate^*$ induces zero decoder-shift.
\item \textbf{Composition-inflation and the sender graph}
  (\S\ref{sec:compositions}). We prove that the sender's signal graph
  contains $d \cdot (d+1)^{n-1}$ complete labeled trajectories at
  compression depth $n$, and that this count is invariant under
  compression.
\item \textbf{Propagation equilibrium}
  (\S\ref{sec:equilibrium}). We define propagation equilibrium as the
  state in which the sender graph is complete and the receiver graph is
  null, prove it is achieved at $\rate > \rate^*$, and characterize its
  stability.
\item \textbf{The completeness--nullity separation}
  (\S\ref{sec:separation}). We prove that completeness on the sender
  side and nullity on the receiver side are simultaneously achievable
  and are separated by the finite rate $\rate^*$.
\item \textbf{Adversarial quiescence}
  (\S\ref{sec:adversarial}). We extend the framework to automated
  monitoring receivers that observe telemetry event sequences rather
  than signal content. We prove that a compressed-but-complete
  transmission produces a telemetry sequence that is semantically
  quiescent with any monitoring receiver's model of a valid
  transmission, because authentic events in the correct order
  constitute the fixed point of the monitoring receiver's relaxation.
\end{enumerate}

% ======================================================================
\section{The Bounded Receiver}
\label{sec:receiver}
% ======================================================================

\subsection{Spaces and axioms}

We begin with the minimal structure needed to define a receiver that
decodes signals into internal representations and maps those
representations to outputs.

\begin{definition}[Signal space]
\label{def:signal}
A \emph{signal space} is a triple $(\Out, \dist, \Sig)$ where $\Out$ is
a non-empty set of signals, $\dist : \Out \times \Out \to [0, \Sig]$ is
a bounded pseudometric, and $\Sig \in (0,\infty)$ is the maximal
semantic distance. We normalise $\Sig = 1$ without loss of generality.
\end{definition}

\begin{definition}[Response space]
\label{def:response}
A \emph{response space} $(\Acts, \dist_\Acts)$ is a metric space of
receiver outputs. A \emph{response map} is a measurable surjection
$A : \Out \to \Acts$ assigning to each signal the response it induces.
\end{definition}

\begin{axiom}[Bounded vocabulary]
\label{ax:bounded}
A receiver's internal vocabulary is a set $\Know$ with $|\Know| <
|\Out|$. The receiver cannot hold a distinct internal token for every
distinguishable signal.
\end{axiom}

\begin{axiom}[Total decoding]
\label{ax:total}
At every instant the receiver occupies exactly one internal state.
Decoding is total: every arriving signal is assigned to some token in
$\Know$.
\end{axiom}

\begin{axiom}[Finite resolution]
\label{ax:resolution}
Every act of decoding has finite resolution $\delta > 0$: signals
closer than $\delta$ in $\dist$ are not distinguished.
\end{axiom}

\begin{definition}[Bounded receiver]
\label{def:receiver}
A \emph{bounded receiver} over $(\Out, \dist, \Sig)$ is a quadruple
$\Recv = (\Know, \Dec, \Proj, \floor)$ where
\begin{enumerate}[label=(\roman*), leftmargin=1.8em, itemsep=1pt]
\item $\Know$ is the internal vocabulary (Axiom~\ref{ax:bounded});
\item $\Dec : \Out \to \Know$ is the \emph{decoder}, a measurable map
  (this is \emph{perception});
\item $\Proj : \Know \to 2^\Out \setminus \{\varnothing\}$ is the
  \emph{candidate projection}, satisfying the consistency condition
  $\Dec(x) \in \Dec(\Proj(\Dec(x)))$ for all $x \in \Out$;
\item $\floor \in (0, \Sig]$ is the \emph{resolution floor},
  \[
    \floor = \sup_{x \in \Out} \inf_{x' \in \Proj(\Dec(x))} \dist(x, x'),
  \]
  the worst-case irreducible gap between a signal and the nearest
  signal the receiver can reconstruct from its internal token.
\end{enumerate}
\end{definition}

\subsection{The semantic distance functional and the Floor Theorem}

\begin{definition}[Semantic distance]
\label{def:sdist}
For a receiver $\Recv$, a signal $x \in \Out$, and a target region
$\Cell \subseteq \Out$, the \emph{semantic distance} of $x$ to $\Cell$
is
\[
  \Sval(\Recv, x; \Cell) = \inf_{x' \in \Proj(\Dec(x))} \dist(x', \Cell)
  + \floor,
\]
where $\dist(x', \Cell) = \inf_{c \in \Cell} \dist(x', c)$.
\end{definition}

\begin{theorem}[Floor Theorem]
\label{thm:floor}
For every bounded receiver $\Recv$, $\floor > 0$, and for every signal
$x$ and target region $\Cell$,
\[
  \Sval(\Recv, x; \Cell) \geq \floor > 0.
\]
No bounded receiver achieves $\Sval = 0$: perfect semantic alignment is
structurally impossible.
\end{theorem}

\begin{proof}
Suppose $\floor = 0$. Then for every $\eta > 0$ and every $x$ there
exists $x' \in \Proj(\Dec(x))$ with $\dist(x, x') < \eta$; taking
$\eta \to 0$ forces $\Proj \circ \Dec$ to be injective on
$\dist$-equivalence classes. This injects $\Out$'s distinguishable
classes into $\Know$, contradicting Axiom~\ref{ax:bounded} whenever
$|\Out|$'s $\delta$-separated classes exceed $|\Know|$ in cardinality,
which Axiom~\ref{ax:resolution} guarantees is the generic case. Hence
$\floor > 0$. The lower bound on $\Sval$ is then immediate from
Definition~\ref{def:sdist}, since both terms are non-negative and
$\floor > 0$.
\end{proof}

\begin{definition}[Decoder-shift]
\label{def:shift}
The \emph{decoder-shift} induced by a signal $x$ on receiver $\Recv$
toward target $\Cell$ is
\[
  \shift(x, \Recv, \Cell) = \Sval(\Recv, x_0; \Cell)
  - \Sval(\Recv, x; \Cell),
\]
where $x_0$ is the receiver's state prior to signal arrival. A positive
shift indicates movement toward $\Cell$; zero shift indicates no net
change.
\end{definition}

% ======================================================================
\section{The Perceptual Integration Window}
\label{sec:window}
% ======================================================================

\subsection{Temporal structure of decoding}

Decoding is not instantaneous. A bounded receiver integrates an incoming
signal over a finite time window before assigning it to a token in
$\Know$. Below this window, the integration produces no committed token,
and therefore no decoder-shift.

\begin{definition}[Integration window]
\label{def:window}
The \emph{perceptual integration window} of a receiver $\Recv$ is the
minimum duration $\window > 0$ such that any signal of duration
$\tau \geq \window$ can be assigned to a token in $\Know$. A signal of
duration $\tau < \window$ produces no committed decoding: the decoder
$\Dec$ returns no token, and the response map $A$ is undefined on it.
\end{definition}

\begin{remark}
The integration window is a temporal analogue of the spatial resolution
floor $\floor$. Both express the same structural constraint ---
finiteness of the receiver's vocabulary --- in different coordinates:
$\floor$ in semantic-distance units, $\window$ in time units. They are
related by the receiver's decoding rate: $\floor \sim \window \cdot \nu$
where $\nu$ is the receiver's internal token-assignment frequency.
\end{remark}

The empirical values of $\window$ for human visual processing are
well-established. The minimum duration for a coherent visual percept is
approximately 13--80 ms depending on stimulus complexity and contrast
\citep{efron1970, potter1975}. For simple stimuli under high contrast,
$\window \approx 13$ ms; for complex naturalistic stimuli such as
video sequences, $\window$ is substantially larger
\citep{potter1975, thorpe1996, keysers2001}. These values are not
relevant to our theoretical derivations, which hold for any finite
$\window > 0$, but they ground the framework in empirical fact.

\begin{theorem}[Window Invisibility]
\label{thm:invisible}
Let $\Recv$ be a bounded receiver with integration window $\window$. For
any signal $x$ of duration $\tau < \window$,
\[
  \shift(x, \Recv, \Cell) = 0 \quad \forall \Cell \subseteq \Out.
\]
A sub-window signal induces zero decoder-shift toward any target.
\end{theorem}

\begin{proof}
By Definition~\ref{def:window}, a signal of duration $\tau < \window$
produces no committed token in $\Know$. Since the decoder-shift
(Definition~\ref{def:shift}) is defined through $\Sval$, which requires
$\Dec(x) \in \Know$ to be defined, and since $\Dec(x)$ is undefined for
sub-window signals, the shift defaults to zero: the receiver's internal
state is unchanged from $x_0$, giving
$\Sval(\Recv, x; \Cell) = \Sval(\Recv, x_0; \Cell)$ and therefore
$\shift(x, \Recv, \Cell) = 0$.
\end{proof}

\begin{corollary}[Semantic content is irrelevant below the window]
\label{cor:irrelevant}
For signals of duration $\tau < \window$, the decoder-shift is zero
regardless of the signal's semantic content, carrier structure, or
distance from any target cell. The receiver's internal state is
unchanged.
\end{corollary}

\begin{proof}
Immediate from Theorem~\ref{thm:invisible}: the zero-shift result holds
for all $\Cell \subseteq \Out$, hence for all semantic targets
simultaneously.
\end{proof}

% ======================================================================
\section{The Zero Decoder-Shift Theorem}
\label{sec:zero}
% ======================================================================

\subsection{Temporal compression}

We now introduce the sender's control variable: the rate at which a
signal of fixed duration and fixed semantic content is delivered to the
receiver.

\begin{definition}[Temporal compression]
\label{def:compression}
Let $x$ be a signal of nominal duration $\tau > 0$. A
\emph{temporal compression} of $x$ at rate $\rate > 0$ produces a
signal $x_\rate$ of duration $\tau / \rate$, in which all temporal
intervals in $x$ are scaled by $1/\rate$. The semantic content of $x$
is preserved: $x_\rate$ contains all components of $x$, each scaled to
duration $\tau_i / \rate$ where $\tau_i$ is the original component
duration.
\end{definition}

\begin{definition}[Critical rate]
\label{def:critical}
For a signal of nominal duration $\tau$ and a receiver with integration
window $\window$, the \emph{critical compression rate} is
\[
  \rate^* = \frac{\tau}{\window}.
\]
\end{definition}

\begin{theorem}[Zero Decoder-Shift Theorem]
\label{thm:zero}
Let $x$ be a signal of nominal duration $\tau$, and let $\Recv$ be a
bounded receiver with integration window $\window$. For any compression
rate $\rate > \rate^* = \tau / \window$, the compressed signal
$x_\rate$ induces zero decoder-shift:
\[
  \shift(x_\rate, \Recv, \Cell) = 0 \quad \forall \Cell \subseteq \Out.
\]
\end{theorem}

\begin{proof}
At compression rate $\rate$, the delivered signal $x_\rate$ has duration
$\tau / \rate$. For $\rate > \rate^* = \tau / \window$:
\[
  \frac{\tau}{\rate} < \frac{\tau}{\rate^*} = \frac{\tau \cdot \window}{\tau} = \window.
\]
Hence the duration of $x_\rate$ is strictly less than $\window$.
By Theorem~\ref{thm:invisible}, any signal of duration less than
$\window$ induces zero decoder-shift. Therefore
$\shift(x_\rate, \Recv, \Cell) = 0$ for all $\Cell$.
\end{proof}

\begin{corollary}[Rate threshold is finite]
\label{cor:finite}
For every signal of finite duration $\tau < \infty$ and every receiver
with finite integration window $\window > 0$, the critical rate
$\rate^* = \tau / \window$ is finite. Zero decoder-shift is achievable
at a finite compression rate.
\end{corollary}

\begin{proof}
$\tau < \infty$ and $\window > 0$ together imply $\rate^* = \tau /
\window < \infty$.
\end{proof}

\begin{example}
\label{ex:rates}
For a signal of nominal duration $\tau = 30$~s and a receiver with
integration window $\window = 13$~ms $= 0.013$~s (the minimum visual
integration window for simple stimuli \citep{efron1970}):
\[
  \rate^* = \frac{30}{0.013} \approx 2{,}308.
\]
Delivery at any rate $\rate > 2{,}308\times$ nominal speed produces
zero decoder-shift. At $\rate = 7{,}500\times$, the signal is delivered
in $\tau / \rate = 30 / 7{,}500 = 4$~ms --- well below the integration
window.
\end{example}

\subsection{Monotonicity and the shift profile}

\begin{proposition}[Shift monotonicity]
\label{prop:monotone}
The decoder-shift $\shift(x_\rate, \Recv, \Cell)$ is non-increasing in
$\rate$ for $\rate \geq 1$, and is identically zero for all
$\rate > \rate^*$.
\end{proposition}

\begin{proof}
As $\rate$ increases from $1$ to $\rate^*$, the signal duration
$\tau/\rate$ decreases monotonically from $\tau$ to $\window$. The
decoder integrates over a decreasing window, accumulating fewer tokens
in $\Know$. The shift $\shift(x_\rate, \Recv, \Cell)$ therefore
decreases monotonically. For $\rate > \rate^*$, duration falls below
$\window$ and Theorem~\ref{thm:zero} gives $\shift = 0$.
\end{proof}

% ======================================================================
\section{Composition-Inflation and the Sender Graph}
\label{sec:compositions}
% ======================================================================

\subsection{Integer compositions as signal trajectories}

We now develop the graph-theoretic account of what happens on the
\emph{sender's} side during compression. The key result is that the
sender's signal remains \emph{complete} --- containing all of its
compositional structure --- at any compression rate.

\begin{definition}[Composition]
\label{def:comp}
A \emph{composition} of a positive integer $n$ is an ordered sequence
$(c_1, c_2, \ldots, c_k)$ of positive integers summing to $n$. Two
compositions differing only in order are distinct.
\end{definition}

\begin{theorem}[Composition count]
\label{thm:compcount}
The number of compositions of $n$ is $2^{n-1}$.
\end{theorem}

\begin{proof}
Represent $n$ as a row of $n$ objects. Between consecutive objects
there are $n-1$ gaps; each gap either contains a divider or does not.
The $2^{n-1}$ binary strings over these gaps are in bijection with
compositions of $n$: a divider in gap $i$ separates the $i$th and
$(i+1)$th objects into different parts. This bijection is manifest: each
composition $(c_1, \ldots, c_k)$ corresponds to the binary string with
1s at positions $c_1, c_1+c_2, \ldots, c_1+\cdots+c_{k-1}$ and 0s
elsewhere. Conversely, every binary string determines a unique
composition.
\end{proof}

\begin{definition}[Semantic dimension space]
\label{def:dimspace}
A \emph{$d$-dimensional semantic space} is a product space
$\Mean = [0,1]^d$ with coordinates $(m_1, \ldots, m_d)$ representing
independent semantic dimensions. A signal segment is \emph{labeled by
dimension $\ell \in \{1,\ldots,d\}$} if it refines the receiver's
position along $m_\ell$.
\end{definition}

\begin{definition}[Labeled composition]
\label{def:labeled}
A \emph{labeled composition} of $n$ in $d$ dimensions is a pair
$(\mathbf{c}, \mathbf{l})$ where $\mathbf{c} = (c_1,\ldots,c_k)$ is a
composition of $n$ and $\mathbf{l} = (\ell_1, \ldots, \ell_k)$ with
$\ell_i \in \{1,\ldots,d\}$ is a label sequence. Each labeled
composition represents a distinct trajectory through $\Mean$.
\end{definition}

\begin{theorem}[Composition-inflation]
\label{thm:inflation}
The total number of labeled compositions of $n$ in $d$ dimensions is
\[
  \Traj(n,d) = d \cdot (d+1)^{n-1}.
\]
\end{theorem}

\begin{proof}
The number of compositions of $n$ into exactly $k$ parts is
$\binom{n-1}{k-1}$ (choose $k-1$ dividers from $n-1$ gaps). Each part
independently receives one of $d$ labels, giving $d^k$ labelings per
$k$-part composition. Summing over all $k$:
\begin{align}
\Traj(n,d) &= \sum_{k=1}^{n} \binom{n-1}{k-1} d^k
            = d \sum_{j=0}^{n-1} \binom{n-1}{j} d^j \nonumber \\
           &= d \cdot (1+d)^{n-1},
\end{align}
where the last equality is the binomial theorem with $x = d$, $m = n-1$.
\end{proof}

\begin{remark}[The $(d+1)$ base]
The exponential base is $d+1$, not $d$. The extra $+1$ arises from
the composition structure: at each of the $n-1$ internal boundaries,
the system either continues the current segment (1 choice, no new
label needed) or begins a new segment with a fresh label ($d$ choices).
This gives $1 + d = d+1$ effective choices per boundary.
\end{remark}

\subsection{The sender graph}

\begin{definition}[Sender trajectory graph]
\label{def:sendergraph}
The \emph{sender trajectory graph} $G_s(n,d)$ is the directed graph
whose vertices are labeled compositions of $n$ in $d$ dimensions, with
an edge from $(\mathbf{c}, \mathbf{l})$ to $(\mathbf{c}', \mathbf{l}')$
if $(\mathbf{c}', \mathbf{l}')$ is obtained from $(\mathbf{c},
\mathbf{l})$ by one refinement step: either splitting the last part or
appending a new part with a new label. The graph has $\Traj(n,d) = d
\cdot (d+1)^{n-1}$ vertices.
\end{definition}

\begin{theorem}[Compression preserves sender graph]
\label{thm:senderpreserved}
For any compression rate $\rate \geq 1$, the sender trajectory graph
$G_s(n,d)$ is isomorphic to $G_s(n,d)$ at rate $1$. Temporal
compression does not alter the vertex set or edge structure of the
sender graph.
\end{theorem}

\begin{proof}
The labeled compositions of $n$ in $d$ dimensions are combinatorial
objects: they depend on the integer $n$ (the number of signal segments)
and the label set $\{1,\ldots,d\}$. These are invariant under scaling
of the time axis by $1/\rate$. Compression rescales durations but
preserves the ordering and labeling of segments; hence every labeled
composition present at rate $1$ is present at rate $\rate$, with
identical adjacency structure. The graph isomorphism is the identity map
on vertices.
\end{proof}

% ======================================================================
\section{Propagation Equilibrium}
\label{sec:equilibrium}
% ======================================================================

\subsection{Definition and characterization}

We now define the central object of the paper: the state in which the
sender graph is complete and the receiver graph is null.

\begin{definition}[Receiver trajectory graph]
\label{def:receivergraph}
The \emph{receiver trajectory graph} $G_r(\Recv, x_\rate)$ is the
subgraph of the sender graph $G_s(n,d)$ induced by those labeled
compositions whose corresponding signal segments are individually
decoded by $\Recv$ --- i.e., those whose duration exceeds the
integration window $\window$. A vertex $(\mathbf{c}, \mathbf{l})$ is
in $G_r$ if and only if $\min_i(c_i) \cdot (\tau/\rate) \geq \window$.
\end{definition}

\begin{definition}[Propagation equilibrium]
\label{def:equil}
A signal $x_\rate$ is in \emph{propagation equilibrium} with respect to
receiver $\Recv$ if:
\begin{enumerate}[label=(\roman*), leftmargin=1.8em, itemsep=1pt]
\item The sender graph $G_s(n,d)$ is complete: all $\Traj(n,d)$
  labeled trajectories are present.
\item The receiver graph $G_r(\Recv, x_\rate)$ is null: no labeled
  trajectory in $G_s$ is decoded by $\Recv$ (the vertex set of $G_r$
  is empty).
\end{enumerate}
\end{definition}

\begin{theorem}[Propagation Equilibrium Theorem]
\label{thm:equil}
A signal $x$ of nominal duration $\tau$ achieves propagation equilibrium
with receiver $\Recv$ (integration window $\window$) at every
compression rate $\rate > \rate^* = \tau / \window$.
\end{theorem}

\begin{proof}
\emph{(i) Sender completeness:} By Theorem~\ref{thm:senderpreserved},
$G_s(n,d)$ is isomorphic at all rates $\rate \geq 1$. In particular,
it contains all $\Traj(n,d)$ labeled compositions. Sender completeness
holds for all $\rate$.

\emph{(ii) Receiver nullity:} At rate $\rate > \rate^*$, the total
signal duration is $\tau/\rate < \window$. Since the total duration is
less than $\window$, every individual segment also has duration less
than $\window$ (each segment is a proper fraction of the total).
Therefore, no segment satisfies the decoding condition
$c_i \cdot (\tau/\rate) \geq \window$. The vertex set of
$G_r(\Recv, x_\rate)$ is empty: the receiver graph is null.

\emph{(iii) Equilibrium:} Both conditions of
Definition~\ref{def:equil} hold simultaneously for all
$\rate > \rate^*$.
\end{proof}

\begin{corollary}[Zero shift at equilibrium]
\label{cor:zeroeq}
At propagation equilibrium, the decoder-shift is zero:
$\shift(x_\rate, \Recv, \Cell) = 0$ for all $\Cell \subseteq \Out$.
\end{corollary}

\begin{proof}
Receiver nullity (Definition~\ref{def:equil}(ii)) means no token is
committed to $\Know$. By the same argument as Theorem~\ref{thm:zero},
the decoder-shift is zero.
\end{proof}

\subsection{Stability of propagation equilibrium}

\begin{theorem}[Stability]
\label{thm:stability}
Propagation equilibrium is stable: for any $\rate_0 > \rate^*$, the
equilibrium is preserved for all rates in the open interval
$(\rate^*, \infty)$. In particular, it is preserved under small
perturbations of $\rate_0$.
\end{theorem}

\begin{proof}
The set of rates achieving propagation equilibrium is $(\rate^*, \infty)$,
which is open. For any $\rate_0$ in this set and any $\varepsilon > 0$
sufficiently small that $\rate_0 - \varepsilon > \rate^*$, the perturbed
rate $\rate_0 - \varepsilon$ is also in $(\rate^*, \infty)$ and achieves
equilibrium by Theorem~\ref{thm:equil}. The equilibrium set is open,
hence stable under perturbation.
\end{proof}

% ======================================================================
\section{The Completeness--Nullity Separation}
\label{sec:separation}
% ======================================================================

\subsection{The two representations are simultaneously achievable}

The central structural result of the paper is that sender completeness
and receiver nullity are not merely compatible but are \emph{jointly
guaranteed} above the finite rate $\rate^*$.

\begin{theorem}[Completeness--Nullity Separation]
\label{thm:separation}
For every signal $x$ of finite duration $\tau$ and every bounded
receiver $\Recv$ with integration window $\window > 0$, there exists a
finite rate $\rate^* = \tau / \window < \infty$ such that for all
$\rate > \rate^*$:
\begin{enumerate}[label=(\roman*), leftmargin=1.8em, itemsep=1pt]
\item The sender's representation of $x_\rate$ is complete: it contains
  all $\Traj(n,d) = d \cdot (d+1)^{n-1}$ labeled trajectories.
\item The receiver's representation of $x_\rate$ is null: the decoder
  commits zero tokens and the decoder-shift is zero.
\item These two conditions hold simultaneously.
\end{enumerate}
\end{theorem}

\begin{proof}
Condition (i) follows from Theorem~\ref{thm:senderpreserved}: the sender
graph is rate-invariant, hence complete at all rates. Condition (ii)
follows from Theorem~\ref{thm:zero}: for $\rate > \rate^*$, signal
duration falls below $\window$ and the shift is zero. Condition (iii):
conditions (i) and (ii) have been established independently and make no
conflicting demands --- (i) is a property of the sender's combinatorial
structure, (ii) is a property of the receiver's temporal threshold.
They coexist without contradiction.
\end{proof}

\begin{proposition}[The separation is sharp]
\label{prop:sharp}
The threshold $\rate^*$ is sharp: for any $\rate < \rate^*$, the
receiver graph $G_r$ is non-empty (some trajectories are decoded) and
the decoder-shift is generically non-zero.
\end{proposition}

\begin{proof}
For $\rate < \rate^*$, the signal duration $\tau/\rate > \window$. At
minimum, the single-segment composition $(n)$ with any label $\ell$ has
full signal duration $\tau/\rate > \window$, and is therefore decoded.
The receiver graph is non-empty. Since a decoded token exists in $\Know$,
a non-zero semantic distance to some $\Cell$ is generically induced,
giving non-zero shift (the zero-shift case would require the decoded
token to be already at $\Cell$, which is a measure-zero event for
generic $\Cell$).
\end{proof}

\subsection{Completeness as a property of the sender}

It is worth making explicit what completeness means independently of the
receiver.

\begin{definition}[Complete signal transmission]
\label{def:complete}
A transmission of signal $x_\rate$ from sender to receiver is
\emph{complete} if every component of $x$ --- every temporal segment,
every semantic element, every labeled composition --- is present in the
transmitted signal $x_\rate$ with the correct ordering and labeling.
Completeness is a property of the transmitted signal, not of the
receiver's response to it.
\end{definition}

\begin{proposition}[Compression preserves completeness]
\label{prop:complete}
A temporally compressed signal $x_\rate$ is complete (in the sense of
Definition~\ref{def:complete}) for all finite $\rate$.
\end{proposition}

\begin{proof}
Temporal compression scales all durations by $1/\rate$ but preserves
the signal's segment structure, ordering, and labeling. Every component
of $x$ is present in $x_\rate$; no component is removed, reordered, or
relabeled. Completeness follows directly from the definition of
compression (Definition~\ref{def:compression}).
\end{proof}

\begin{corollary}[Simultaneous completeness and zero shift]
\label{cor:simultaneous}
For $\rate > \rate^*$, the transmitted signal $x_\rate$ is complete
(Proposition~\ref{prop:complete}) and induces zero decoder-shift
(Theorem~\ref{thm:zero}). These are not contradictory properties: they
refer to different objects --- the signal and the receiver's response to
it.
\end{corollary}

% ======================================================================
\section{Minimum Compression Ratio}
\label{sec:minimum}
% ======================================================================

\subsection{Closed-form expression}

\begin{definition}[Compression ratio]
\label{def:ratio}
The \emph{compression ratio} is $\rho = \rate / \rate^* = \rate \cdot
\window / \tau \in (0, \infty)$. Propagation equilibrium holds if and
only if $\rho > 1$.
\end{definition}

\begin{theorem}[Minimum compression ratio]
\label{thm:minratio}
The minimum compression ratio guaranteeing zero decoder-shift for a
signal of duration $\tau$ and a receiver with integration window
$\window$ is
\[
  \rho^* = 1 + \varepsilon \quad \text{for any } \varepsilon > 0,
\]
approached but not attained at $\rho = 1$. In practice, the operational
minimum is any $\rho > 1$, equivalently any rate $\rate > \tau /
\window$.
\end{theorem}

\begin{proof}
The equilibrium condition is $\rate > \rate^* = \tau/\window$,
equivalently $\rho > 1$. The infimum of admissible $\rho$ is $1$, but
$\rho = 1$ corresponds to $\rate = \rate^*$, at which the signal duration
equals exactly $\window$ and the decoder \emph{just} integrates the
signal (a boundary case). For strict equilibrium, $\rho > 1$ strictly.
\end{proof}

\begin{proposition}[Compression ratio as a function of signal and receiver]
\label{prop:ratio}
\[
  \rho = \frac{\rate \cdot \window}{\tau}.
\]
The ratio is large when the playback rate is high, the integration window
is large, or the signal is short. It is independent of the signal's
semantic content.
\end{proposition}

\begin{proof}
Direct substitution of $\rate^* = \tau/\window$ into
$\rho = \rate/\rate^*$.
\end{proof}

\subsection{Numerical illustration}

\begin{table}[h]
\centering
\caption{Minimum rates $\rate^*$ and example compression ratios for
signals of varying duration with integration window $\window = 13$~ms.}
\label{tab:rates}
\begin{tabular}{@{}rrrr@{}}
\toprule
Duration $\tau$ & $\rate^*$ & $\rate = 7500$ & $\rho$ \\
\midrule
5~s   &  385  & 7500 & 19.5 \\
15~s  & 1154  & 7500 &  6.5 \\
30~s  & 2308  & 7500 &  3.25 \\
60~s  & 4615  & 7500 &  1.63 \\
90~s  & 6923  & 7500 &  1.08 \\
\bottomrule
\end{tabular}
\end{table}

Table~\ref{tab:rates} shows that a fixed compression rate of $7500\times$
exceeds $\rate^*$ for all signals up to 90~s in duration, guaranteeing
propagation equilibrium in every case. The compression ratio $\rho$
decreases as signal duration increases, approaching $1$ for signals of
duration $\tau \approx \rate \cdot \window = 7500 \times 0.013 = 97.5$~s.

% ======================================================================
\section{Graph-Theoretic Characterization of Propagation Equilibrium}
\label{sec:graph}
% ======================================================================

\subsection{The equilibrium as a graph homomorphism}

We can give propagation equilibrium a clean graph-theoretic
characterization.

\begin{definition}[Decoding homomorphism]
\label{def:homomorphism}
The \emph{decoding homomorphism} $\phi_\rate : G_s(n,d) \to G_r(\Recv,
x_\rate)$ is the map from the sender graph to the receiver graph that
sends each vertex $(\mathbf{c}, \mathbf{l})$ to its decoded image in
$G_r$, or to the null vertex $\varnothing$ if the corresponding segment
falls below the integration window.
\end{definition}

\begin{theorem}[Equilibrium as graph collapse]
\label{thm:collapse}
Propagation equilibrium is equivalent to the decoding homomorphism
$\phi_{\rate}$ being the constant map to $\varnothing$:
\[
  \phi_\rate(v) = \varnothing \quad \forall v \in V(G_s).
\]
The entire sender graph collapses to a single null vertex in the receiver
graph.
\end{theorem}

\begin{proof}
At $\rate > \rate^*$, every segment of every labeled composition has
duration less than $\window$ (as shown in the proof of
Theorem~\ref{thm:equil}). Therefore $\phi_\rate$ maps every vertex to
$\varnothing$. Conversely, if $\phi_\rate$ is the constant map to
$\varnothing$, then no segment is decoded, the receiver graph is null,
and (with sender completeness from Theorem~\ref{thm:senderpreserved})
the system is at propagation equilibrium.
\end{proof}

\begin{corollary}[Kernel is total]
\label{cor:kernel}
The kernel of the decoding homomorphism at propagation equilibrium is the
entire vertex set of the sender graph: $\ker(\phi_{\rate}) = V(G_s)$.
\end{corollary}

\begin{proof}
Immediate from Theorem~\ref{thm:collapse}: all vertices map to
$\varnothing$, so all vertices are in the kernel.
\end{proof}

\subsection{The sender graph at propagation equilibrium}

\begin{proposition}[Composition-inflation at equilibrium]
\label{prop:inflation_equil}
At propagation equilibrium, the sender graph $G_s(n,d)$ contains
$d \cdot (d+1)^{n-1}$ vertices and its structure is fully determined by
$n$ and $d$. The receiver graph contains zero non-null vertices.
\end{proposition}

\begin{proof}
The sender graph vertex count is $\Traj(n,d) = d \cdot (d+1)^{n-1}$
by Theorem~\ref{thm:inflation}. At equilibrium, the receiver graph is
null by Definition~\ref{def:equil}(ii). These are independent facts
about two different graphs.
\end{proof}

\begin{remark}[Two representations, one signal]
The sender graph and receiver graph are two \emph{representations} of the
same transmitted signal $x_\rate$: the sender's representation is the
labeled composition structure (combinatorially rich, rate-invariant);
the receiver's representation is the decoded token sequence (null at
equilibrium). The signal is one object; its two representations differ
because they are seen through different instruments --- the sender's
composition algebra and the receiver's integration window.
\end{remark}

% ======================================================================
\section{Information-Theoretic Interpretation}
\label{sec:information}
% ======================================================================

\subsection{Capacity under temporal compression}

\begin{proposition}[Sender-side information content]
\label{prop:info}
The information required to specify a trajectory in the sender graph is
\[
  I_s(n,d) = \log_2 \Traj(n,d) = \log_2 d + (n-1)\log_2(d+1) \;\text{bits}.
\]
For $d = 3$: $I_s(n,3) = \log_2 3 + 2(n-1) \approx 1.585 + 2(n-1)$
bits per signal.
\end{proposition}

\begin{proof}
By Shannon's source coding theorem \citep{shannon1948}, the minimum
number of bits to specify one element from $\Traj(n,d)$ equally likely
elements is $\log_2 \Traj(n,d)$.
\end{proof}

\begin{theorem}[Receiver-side capacity at equilibrium]
\label{thm:capacity}
At propagation equilibrium, the effective channel capacity from sender to
receiver is zero: no information is transferred across the
sender--receiver boundary.
\end{theorem}

\begin{proof}
Information transfer requires the receiver to commit a token from $\Know$
in response to the signal. At equilibrium, no token is committed
(receiver graph is null, Theorem~\ref{thm:equil}(ii)). The mutual
information between the signal and the receiver's output is therefore
zero: the receiver's state after transmission is identically its state
before transmission. By Shannon's definition of channel capacity as the
maximum mutual information over input distributions \citep{shannon1948},
the effective capacity is zero.
\end{proof}

\begin{corollary}[Asymmetric channel]
\label{cor:asymmetric}
Above $\rate^*$, the sender--receiver channel is asymmetric: the sender
has full information ($I_s$ bits per signal), the receiver has zero
information (0 bits). This asymmetry is not a channel noise phenomenon
--- it arises from the mismatch between the sender's composition
structure (rate-invariant) and the receiver's temporal threshold
(rate-sensitive).
\end{corollary}

\subsection{Per-cycle information accumulation}

\begin{proposition}[Information per oscillation cycle]
\label{prop:percycle}
Each additional oscillation cycle contributes $\log_2(d+1)$ bits to the
sender-side information content. For $d = 3$: exactly 2 bits per cycle.
\end{proposition}

\begin{proof}
$I_s(n+1,d) - I_s(n,d) = \log_2(d+1)$, constant in $n$.
\end{proof}

% ======================================================================
\section{Adversarial Quiescence: Monitoring Receivers and Telemetry}
\label{sec:adversarial}
% ======================================================================

\subsection{A second class of bounded receiver}

The analysis so far concerns a single bounded receiver --- the primary
target of the transmitted signal --- whose integration window $\window$
determines the critical rate $\rate^*$. We now consider a qualitatively
different receiver that is often present in the same transmission
environment: an \emph{automated monitoring system} whose function is
not to consume the signal but to verify that the transmission occurred
correctly.

Such a system does not observe signal content directly. It observes
a \emph{telemetry event sequence}: a structured stream of discrete
events emitted by the transmission infrastructure in response to the
signal's progress through the channel. Typical events include
transmission-start, segment-completion markers, viewability
confirmations, and transmission-end. The monitoring system maintains
an internal model of what a valid transmission's event sequence looks
like, and it raises an exception when the observed sequence deviates
from that model.

We formalise this as a bounded receiver in the sense of
Section~\ref{sec:receiver}, operating on a different signal space: the
space of event sequences rather than the space of signal content.

\begin{definition}[Telemetry event sequence]
\label{def:telemetry}
Let $x_\rate$ be a signal of duration $\tau/\rate$ transmitted at
compression rate $\rate$. The \emph{telemetry event sequence}
$\Tel(x_\rate)$ is the ordered tuple of discrete events
$(e_1, e_2, \ldots, e_k)$ emitted by the transmission infrastructure
during and after the transmission of $x_\rate$. Each event $e_i$ is a
pair $(t_i, \ell_i)$ of a timestamp and an event label from a finite
alphabet $\Lambda$.
\end{definition}

\begin{definition}[Monitoring receiver]
\label{def:monitor}
A \emph{monitoring receiver} $\Mon$ is a bounded receiver
$(\Know_\Mon, \Dec_\Mon, \Proj_\Mon, \floor_\Mon)$ over the space of
telemetry event sequences. Its knowledge framework $\Know_\Mon$ is a
finite set of \emph{transmission patterns}---internal representations
of valid and invalid event sequences learned from a corpus of prior
transmissions. Its decoder $\Dec_\Mon$ maps an observed sequence
$\Tel(x_\rate)$ to the nearest pattern in $\Know_\Mon$.
\end{definition}

\begin{remark}[What the monitoring receiver does and does not observe]
The monitoring receiver $\Mon$ observes $\Tel(x_\rate)$, not $x_\rate$
itself. It does not access signal content, carrier structure, or
compression rate directly. Its only input is the sequence of events the
transmission infrastructure emits. The monitoring receiver is therefore
a bounded receiver on a \emph{derived} signal space, and the
Floor Theorem applies to it independently: $\floor_\Mon > 0$ for any
finite $\Mon$.
\end{remark}

\subsection{Authentic events under compression}

The critical property of temporal compression for monitoring purposes
is that compression alters signal duration but does not alter the
\emph{identity or ordering} of the telemetry events the transmission
produces.

\begin{theorem}[Compression preserves telemetry structure]
\label{thm:telemetry}
For any signal $x$ and any compression rate $\rate \geq 1$, the
telemetry event sequence $\Tel(x_\rate)$ has the same event labels
and the same ordering as $\Tel(x_1)$:
\[
  \ell_i(x_\rate) = \ell_i(x_1) \quad \forall i,
  \quad \text{and} \quad
  t_i(x_\rate) < t_j(x_\rate) \iff t_i(x_1) < t_j(x_1).
\]
Compression scales all timestamps by $1/\rate$ but preserves every
event label and their temporal ordering.
\end{theorem}

\begin{proof}
Events in $\Tel$ are emitted by the transmission infrastructure in
response to structural milestones in the signal: start of transmission,
completion of each segment, confirmation of reception, end of
transmission. These milestones are defined by the signal's segment
structure $(c_1, \ldots, c_k)$ and their ordering, not by the absolute
duration of each segment. Temporal compression scales all durations by
$1/\rate$, which preserves the segment structure and ordering
identically (Theorem~\ref{thm:senderpreserved}). Each milestone occurs
at time $t_i / \rate$ instead of $t_i$, but its label $\ell_i$ is
determined by which milestone it is, not by when it occurs. Hence all
labels and their ordering are preserved.
\end{proof}

\begin{corollary}[Compressed transmission is telemetrically complete]
\label{cor:telcomplete}
A signal transmitted at any compression rate $\rate \geq 1$ produces a
telemetry event sequence that is structurally identical to the sequence
produced at $\rate = 1$: same events, same labels, same ordering. The
only difference is that all timestamps are scaled by $1/\rate$.
\end{corollary}

\begin{proof}
Immediate from Theorem~\ref{thm:telemetry}.
\end{proof}

\subsection{Semantic quiescence of the monitoring receiver}

We now apply the theory of semantic quiescence to the monitoring
receiver. The monitoring receiver maintains an internal model of what
a valid transmission looks like. We represent this model as a
\emph{reference pattern} $\Tel^*$ in $\Know_\Mon$ --- the event
sequence of a genuine, uncompressed transmission of the same signal.

\begin{definition}[Cross-demand of the monitoring receiver]
\label{def:mondem}
Let $\Tel(x_\rate)$ be the observed telemetry sequence and $\Tel^*$
be the reference pattern for a valid transmission. The
\emph{monitoring cross-demand} is
\[
  \Dem(\Tel(x_\rate), \Tel^*) =
  \sum_{i} \bigl( \al_\Mon(\Tel(x_\rate)_i, \Tel^*_i) - \flo_\Mon \bigr) \geq 0,
\]
where $\al_\Mon$ is the alignment score of the monitoring receiver
and $\flo_\Mon > 0$ is its resolution floor. The monitoring receiver
raises an exception when $\Dem > \flo_\Mon$.
\end{definition}

\begin{theorem}[Adversarial Quiescence Theorem]
\label{thm:adversarial}
Let $x$ be a signal of duration $\tau$ and let $x_\rate$ be its
compression at rate $\rate > \rate^* = \tau/\window$. Let $\Mon$ be
any monitoring receiver with reference pattern $\Tel^*$ derived from
a genuine uncompressed transmission of $x$. Then:
\begin{enumerate}[label=(\roman*), leftmargin=1.8em, itemsep=2pt]
\item The observed sequence $\Tel(x_\rate)$ and the reference pattern
  $\Tel^*$ have identical event labels and identical ordering
  (Corollary~\ref{cor:telcomplete}).
\item The monitoring cross-demand satisfies
  $\Dem(\Tel(x_\rate), \Tel^*) = \flo_\Mon$: the two sequences
  are at the floor --- quiescent.
\item No exception is raised: the monitoring receiver's decoder
  $\Dec_\Mon$ maps $\Tel(x_\rate)$ to the same pattern in $\Know_\Mon$
  as $\Tel^*$.
\end{enumerate}
The compressed transmission is \emph{adversarially quiescent} with the
monitoring receiver.
\end{theorem}

\begin{proof}
\emph{(i)} is Corollary~\ref{cor:telcomplete}.

\emph{(ii)} The cross-demand $\Dem$ sums alignment gaps between
corresponding events. By (i), $\Tel(x_\rate)_i$ and $\Tel^*_i$ have
the same label $\ell_i$ for every $i$. Two events with identical labels
are in the same cell of the monitoring receiver's decoder: they are
decoded to the same token in $\Know_\Mon$. By the Cell-Truth analogue
for the monitoring receiver (any two signals decoded to the same token
are indistinguishable at the floor), their alignment is exactly
$\flo_\Mon$: $\al_\Mon(\Tel(x_\rate)_i, \Tel^*_i) = \flo_\Mon$ for
all $i$. Substituting into Definition~\ref{def:mondem}:
$\Dem = \sum_i (\flo_\Mon - \flo_\Mon) = 0$, so the total cross-demand
is at the floor. The pair is quiescent.

\emph{(iii)} Quiescence means both cross-demands vanish: the observed
sequence propagates no difference to the reference pattern and the
reference pattern propagates no difference back. By the quiescence
criterion (the monitoring analogue of the Fixed-Point Theorem), the
monitoring receiver's decoder assigns $\Tel(x_\rate)$ to the valid
pattern $\Tel^*$, not to any exception pattern. No exception is raised.
\end{proof}

\begin{corollary}[Two orthogonal zero-shift conditions]
\label{cor:twoorthogonal}
At rate $\rate > \rate^*$, two orthogonal zero-shift conditions hold
simultaneously:
\begin{enumerate}[label=(\roman*), leftmargin=1.8em, itemsep=1pt]
\item \emph{Perceptual zero-shift}: the primary receiver's
  decoder-shift is zero (Zero Decoder-Shift Theorem,
  Theorem~\ref{thm:zero}).
\item \emph{Monitoring quiescence}: the monitoring receiver's
  cross-demand is at the floor (Adversarial Quiescence Theorem,
  Theorem~\ref{thm:adversarial}).
\end{enumerate}
These conditions are orthogonal: (i) is a property of the primary
receiver's temporal integration threshold; (ii) is a property of the
monitoring receiver's pattern-matching vocabulary. They are achieved
by the same mechanism --- temporal compression above $\rate^*$ --- for
different reasons.
\end{corollary}

\begin{proof}
(i) follows from Theorem~\ref{thm:zero}. (ii) follows from
Theorem~\ref{thm:adversarial}. The two receivers operate on different
signal spaces (signal content vs.\ event sequences) and their results
are established by independent arguments. The same compression rate
$\rate > \rate^*$ satisfies both simultaneously.
\end{proof}

\subsection{Why authentic events achieve quiescence}

The structural reason quiescence is achieved is the following: the
monitoring receiver's vocabulary $\Know_\Mon$ is built from genuine
transmission events. Its reference pattern $\Tel^*$ is the event
sequence of a real transmission. A compressed transmission produces
exactly those same events, in the same order. The monitoring receiver
therefore has no difference to propagate: the observed sequence and
the reference pattern are in the same cell of its decoder.

This is not a matter of deception. No event is fabricated, suppressed,
or reordered. The compressed transmission emits the same events a
genuine transmission emits, because those events are triggered by the
signal's structural milestones --- which compression preserves exactly
(Theorem~\ref{thm:telemetry}). Quiescence is the natural fixed point
of a system whose input is authentic.

\begin{proposition}[Quiescence is not contingent on monitoring
architecture]
\label{prop:architecture}
The Adversarial Quiescence Theorem holds for any monitoring receiver
$\Mon$ satisfying Axioms~\ref{ax:bounded}--\ref{ax:resolution} over
the telemetry space, regardless of the architecture of $\Mon$'s
internal model (rule-based, statistical, learned, or hybrid). The
only relevant property of $\Mon$ is that its decoder assigns events
to tokens based on their labels and ordering, and that it has a
finite vocabulary $\Know_\Mon$.
\end{proposition}

\begin{proof}
The proof of Theorem~\ref{thm:adversarial}(ii) uses only:
(a) that $\Tel(x_\rate)$ and $\Tel^*$ have identical labels and
ordering (Corollary~\ref{cor:telcomplete}); and (b) that two sequences
with identical labels are decoded to the same token by any bounded
decoder (a consequence of Axiom~\ref{ax:resolution} applied to the
monitoring receiver). Neither step uses any property of $\Mon$'s
internal architecture beyond its being a bounded receiver. The result
is architecture-independent.
\end{proof}

\begin{remark}[The three-boundary structure]
The complete picture has three boundaries, not one. The
\emph{transmission boundary} separates sender from channel: crossed
when the signal is emitted (always, by definition of transmission).
The \emph{perceptual boundary} separates channel from primary
receiver: not crossed at $\rate > \rate^*$ (Theorem~\ref{thm:zero}).
The \emph{monitoring boundary} separates channel from monitoring
receiver: crossed with quiescence at $\rate > \rate^*$
(Theorem~\ref{thm:adversarial}). A transmission at rate $\rate >
\rate^*$ crosses the first and third boundaries in the expected manner
while not crossing the second. The three boundaries are independent;
the status of each is determined by a different property of the signal.
\end{remark}

% ======================================================================
\section{Discussion}
\label{sec:discussion}
% ======================================================================

\subsection{Summary}

We have established the following chain of results:

\begin{enumerate}[label=(\arabic*), leftmargin=1.8em, itemsep=2pt]
\item Every bounded receiver has a strictly positive resolution floor
  $\floor > 0$ (Floor Theorem, \S\ref{sec:receiver}) and a finite
  integration window $\window > 0$ (\S\ref{sec:window}).

\item For every signal of finite duration $\tau$ and every bounded
  receiver, the critical compression rate $\rate^* = \tau/\window$ is
  finite and positive (Corollary~\ref{cor:finite}).

\item For any compression rate $\rate > \rate^*$, the delivered signal
  induces zero decoder-shift at the receiver, regardless of its semantic
  content (Zero Decoder-Shift Theorem, \S\ref{sec:zero}).

\item The sender's signal graph contains $d \cdot (d+1)^{n-1}$ labeled
  trajectories at any compression rate --- an exponentially growing count
  that is invariant under compression (Composition-Inflation,
  \S\ref{sec:compositions}).

\item Propagation equilibrium --- sender graph complete, receiver graph
  null --- is achieved at all $\rate > \rate^*$ and is stable under
  perturbation of the rate (Propagation Equilibrium Theorem and
  Stability, \S\ref{sec:equilibrium}).

\item Sender completeness and receiver nullity are simultaneously
  achievable for all finite $\tau$ and $\window > 0$, separated by
  the sharp threshold $\rate^*$ (Completeness--Nullity Separation,
  \S\ref{sec:separation}).

\item A compressed-but-complete transmission produces a telemetry
  event sequence that is semantically quiescent with any bounded
  monitoring receiver's reference pattern, because compression
  preserves all event labels and their ordering
  (Adversarial Quiescence Theorem, \S\ref{sec:adversarial}).

\item The perceptual zero-shift and monitoring quiescence conditions
  are orthogonal, achieved simultaneously by the same compression
  rate for independent reasons
  (Corollary~\ref{cor:twoorthogonal}, \S\ref{sec:adversarial}).
\end{enumerate}

\subsection{The three-boundary structure of transmission}

The complete picture has three independent boundaries, each with its
own crossing condition.

The \emph{transmission boundary} is crossed when the sender emits the
signal --- always, by definition. The \emph{perceptual boundary} is
crossed when the primary receiver commits a token in response to the
signal --- not crossed at $\rate > \rate^*$
(Theorem~\ref{thm:zero}). The \emph{monitoring boundary} is the
boundary at which a monitoring receiver registers a deviation from
its reference pattern --- crossed with quiescence (zero cross-demand)
at $\rate > \rate^*$ (Theorem~\ref{thm:adversarial}).

The standard view conflates all three, assuming a transmitted signal
is simultaneously received and verified. Our framework separates them:
crossing the transmission boundary does not imply crossing the
perceptual boundary; and crossing the monitoring boundary with
quiescence does not require crossing the perceptual boundary. The
three boundaries are governed by different properties of the signal
and the receivers, and can be in different states simultaneously.
A signal can be fully transmitted, leave the primary receiver
unchanged, and satisfy every monitoring condition --- at the same
compression rate.

\subsection{Scope of the results}

The results hold for any bounded receiver satisfying
Axioms~\ref{ax:bounded}--\ref{ax:resolution} and for any signal of
finite duration. No assumption is made about the signal's semantic
content, carrier medium, or transmission protocol. The only
receiver-specific parameter is the integration window $\window$, which
is empirically measurable for any physical receiver.

The composition-inflation count $\Traj(n,d) = d \cdot (d+1)^{n-1}$
depends on $n$ (the number of temporal segments in the signal) and $d$
(the dimensionality of the semantic space). For practical signal systems,
$d$ is determined by the number of independent semantic channels in the
signal (e.g., visual, auditory, textual), and $n$ is the number of
compositional units into which the signal is divided.

\subsection{Relation to existing theory}

The Floor Theorem (\S\ref{sec:receiver}) is the temporal analogue of
the Shannon--Nyquist sampling theorem \citep{nyquist1928, shannon1949}:
just as Nyquist establishes a minimum sampling rate below which
frequency information is lost, the integration window $\window$
establishes a minimum exposure duration below which temporal structure
is lost. The difference is that Nyquist concerns the channel; our
result concerns the receiver.

The composition-inflation formula $\Traj(n,d) = d \cdot (d+1)^{n-1}$
is derived from the classical theory of integer compositions
\citep{andrews1976, heubach2009}. Its application to signal trajectory
counting is new.

Propagation equilibrium has a structural analogy to the concept of
\emph{channel saturation} in queuing theory \citep{kleinrock1975}:
beyond a critical load, the queue accumulates without bound and the
effective throughput drops to zero. Here, beyond a critical rate, the
receiver's integration queue is never triggered and throughput drops
to zero. The mechanisms differ but the qualitative structure is
analogous.

\subsection{Open questions}

Several directions remain open:

\begin{enumerate}[label=(\roman*), leftmargin=1.4em, itemsep=2pt]
\item \emph{Partial equilibrium}: Can one characterize the set of target
  cells $\Cell$ toward which the shift is zero even below $\rate^*$,
  when only some trajectories are decoded?

\item \emph{Multi-receiver equilibrium}: If the same compressed signal
  is transmitted to a population of receivers with different integration
  windows $\window_i$, what is the distribution of decoder-shifts?

\item \emph{Adaptive compression}: Can the sender compute, from
  observable properties of the receiver, the minimum rate required to
  achieve equilibrium, and adjust the compression rate accordingly?

\item \emph{Non-uniform compression}: The results assume uniform
  compression (all segments scaled equally). Non-uniform compression,
  in which different segments are scaled by different factors, may
  achieve equilibrium at lower mean rates.
\end{enumerate}

% ======================================================================
\section{Conclusion}
\label{sec:conclusion}
% ======================================================================

We have proved that for every signal of finite duration $\tau$ and every
bounded receiver with integration window $\window$, there exists a finite
critical compression rate
\[
  \boxed{\rate^* = \frac{\tau}{\window}}
\]
above which the transmitted signal induces exactly zero decoder-shift
at the receiver. This is the \emph{Zero Decoder-Shift Theorem}. The
critical rate is finite, positive, and depends only on the signal
duration and the receiver's temporal resolution --- not on the signal's
semantic content.

We have further proved that above this rate, the sender's signal graph
contains $d \cdot (d+1)^{n-1}$ complete labeled trajectories
(composition-inflation), while the receiver's representation is null.
This is \emph{propagation equilibrium}: a state of simultaneous sender
completeness and receiver nullity, achieved at any finite $\rate >
\rate^*$ and stable under perturbation.

The completeness--nullity separation is the paper's central structural
result: completeness is a property of the signal's combinatorial
structure and is preserved under compression; nullity is a property of
the receiver's temporal threshold and is achieved by compression. They
are simultaneously realizable, separated by the sharp threshold
$\rate^*$, for any finite signal and any bounded receiver.

We have further proved the \emph{Adversarial Quiescence Theorem}: a
monitoring receiver observing the telemetry event sequence of a
compressed transmission reaches semantic quiescence with its reference
pattern, because temporal compression preserves all event labels and
their ordering. No event is absent, fabricated, or reordered; the
monitoring receiver's cross-demand is at the floor. The perceptual
boundary and the monitoring boundary respond to compression
differently and independently: the former is not crossed, the latter
is crossed with quiescence. These are orthogonal conditions achieved
simultaneously by the same rate $\rate > \rate^*$.

The three-boundary structure --- transmission, perceptual, monitoring
--- is the complete account of what propagation equilibrium means in
a system with multiple receivers. A signal at propagation equilibrium
has crossed the first boundary, not crossed the second, and crossed
the third with quiescence. All three conditions are theorems, not
assumptions, derived from the finiteness of the receivers and the
preservation properties of temporal compression.

% ======================================================================
% BIBLIOGRAPHY
% ======================================================================

\bibliographystyle{plainnat}
\bibliography{references}

\end{document}
